\documentclass[12pt]{article}
\usepackage[utf8]{inputenc}
\usepackage[T1]{fontenc}
\usepackage{hyperref}
\title{Western Punjabi Language Resources for K-12 Curriculum Development}
\author{Compiled by Claude Sonnet 5.5 (claude-sonnet-5-5) for Andr\'as Kornai}
\date{}
\begin{document}
\maketitle
\section{Introduction}
Western Punjabi (ISO 639-3 pnb) is an Indo-Aryan language of Pakistan and India, part of the Lahnda grouping and written in the Perso-Arabic Shahmukhi script. Searches returned mostly historical works; no modern textbooks were found.
\section{Primary Educational Resources}
No modern Shahmukhi textbooks or school programmes were found in searches.
\section{Linguistic and Reference Materials}
\begin{itemize}
\item Lahnda, overview of the varieties and terminology: \url{https://en.wikipedia.org/wiki/Lahnda}
\item Grammar and Dictionary of Western Panjabi as Spoken in the Shahpur District (J. Wilson, 1899), scanned text: \url{https://dn790001.ca.archive.org/0/items/WesternPunjabiShahpurDistrict/Western%20Punjabi%20%28Shahpur%20District%29_text.pdf}
\item The same work, rare-books record: \url{https://icvtesting.nvli.in/rarebooks/grammar-and-dictionary-western-panjabi-spoken-shahpur-district-proverbs-sayings-and-0}
\item Dictionary of the Jatki or Western Panjabi language (Jukes, 1900): \url{https://search.worldcat.org/oclc/1128095742}
\item Routledge reprint of a Western Punjabi reference work: \url{https://www.taylorfrancis.com/books/9780367248673}
\end{itemize}
\section{Religious and Community Resources}
None found in searches.
\section{Digital and Audio Resources}
None found in searches.
\section{Curriculum Development Context}
The language has a very large speaker base but the materials found are century-old. A curriculum would need modern, Shahmukhi-based materials and should clarify the relation to standard Punjabi.
\end{document}